\documentclass[11pt,a4paper]{article}
\usepackage{turkos}

\begin{document}
\phaseheader{7}{Kernel Heap and Core Data Structures}{Give the kernel malloc and free, a reusable linked list, and an automated test run that guards everything you have built so far}{42}{48}

\ataglance{1--2 weeks}{3}{Phase 6: paging, the direct map, the VMM, pre-allocated kernel page tables}{Whole-page allocation
through the direct map, a growable kernel heap with \code{kmalloc}, \code{kfree}, \code{kcalloc} and \code{krealloc}
(first fit, splitting, coalescing, corruption checks), an intrusive doubly linked list, and \code{make test}, which
boots a headless QEMU, runs every test and exits with pass or fail.}

\section{Why this phase}

The next phases create objects whose number and size you can't know in advance: a task structure for every thread,
a file object for every \code{open}, a buffer for every disk block, a copy of every path name. They need a general
allocator, and the PMM only hands out whole 4\,KiB frames. The kernel heap sits on top of the PMM and the VMM and
cuts pages into objects of any size.

Writing an allocator is also the best practice you can get for the kind of code that fills a kernel: pointer
arithmetic, invariants that must hold after every operation, and bugs that show up far from where they were made.
That is why this phase also builds the automated test run. From now on, one command boots Turk-OS without a window,
runs every test, and tells you in seconds whether today's change broke something from three phases ago.

\section{Concepts you need}

\subsection{Free-space management}

\OSTEP chapter 17 is the core reading. An allocator manages a region as a sequence of \textbf{blocks}, each with a
small \textbf{header} in front of the memory it hands out. The header records the block's size and whether it is
free, so \code{kfree(p)} can find it from the pointer alone by stepping back one header. To allocate, the allocator
searches the free blocks for one that is large enough, according to a \textbf{strategy}: first fit (the first one
that fits), best fit (the smallest that fits), or next fit (first fit starting where the last search stopped). If the
chosen block is much larger than needed, it \textbf{splits} it. When a block is freed next to another free block, it
\textbf{coalesces} them into one, otherwise memory breaks into fragments that are each too small to use
(Figure~\ref{fig:heap}).

\begin{figure}[htbp]
\centering
\begin{tikzpicture}[h/.style={draw=tkNavy!70, fill=tkNavy!15, minimum height=0.7cm, minimum width=0.7cm, font=\sffamily\tiny},
  u/.style={draw=tkNavy!70, fill=tkRed!10, minimum height=0.7cm, font=\sffamily\scriptsize},
  f/.style={draw=tkNavy!70, fill=tkGreen!12, minimum height=0.7cm, font=\sffamily\scriptsize}]
  \node[font=\sffamily\scriptsize\bfseries, anchor=east] at (-0.2,0) {before \code{kfree(B)}};
  \node[h, anchor=west] (h1) at (0,0) {hdr};
  \node[f, minimum width=2.2cm, right=0pt of h1] (a) {A: free};
  \node[h, right=0pt of a] (h2) {hdr};
  \node[u, minimum width=2.6cm, right=0pt of h2] (b) {B: in use};
  \node[h, right=0pt of b] (h3) {hdr};
  \node[f, minimum width=1.8cm, right=0pt of h3] (c) {C: free};
  \node[h, right=0pt of c] (h4) {hdr};
  \node[u, minimum width=2.0cm, right=0pt of h4] (d) {D: in use};
  \node[font=\sffamily\scriptsize\bfseries, anchor=east] at (-0.2,-1.3) {after coalescing};
  \node[h, anchor=west] (g1) at (0,-1.3) {hdr};
  \node[f, minimum width=8.7cm, right=0pt of g1] (abc) {one free block: A + header + B + header + C};
  \node[h, right=0pt of abc] (g4) {hdr};
  \node[u, minimum width=2.0cm, right=0pt of g4] (d2) {D: in use};
  \draw[tkarrow] (b.south) -- ([yshift=-0.3cm]b.south) -- ([yshift=0.1cm]abc.north -| b.south);
\end{tikzpicture}
\caption{Freeing B between two free neighbours merges all three into a single block. Without coalescing, a request larger than A, B or C alone would fail.}
\label{fig:heap}
\end{figure}

\subsection{Kernel allocators in practice}

Real kernels usually combine several allocators, and \OSC\,\S10.8 describes the two classic ones. The \textbf{buddy
system} manages power-of-two blocks that can split in half and merge back cheaply. \textbf{Slab allocation} keeps a
cache of ready-made objects of one type (say, every \code{struct task}), which avoids fragmentation and
re-initialisation. Turk-OS uses a simple layered scheme that is easy to get right and easy to replace later. Whole
pages come straight from the PMM through the direct map: \code{kalloc_page()} is just \code{P2V(pmm_alloc_frame())},
and it gives you page-aligned memory for page tables and similar structures. Everything smaller comes from the heap,
a first-fit list in the virtual region \code{0xE0000000}--\code{0xEFFFFFFF} that grows by mapping new frames as needed.

\subsection{Intrusive lists}

A textbook linked list allocates a separate node that points at the data. Kernels usually do it the other way round:
the list links live \emph{inside} the object (\code{struct task} contains a \code{struct list_node}), and a macro,
\code{container_of}, gets from a pointer to the embedded node back to the object that contains it. The list code never
allocates memory, an object can sit on several lists at once, and a single generic implementation serves every list
in the kernel. The same design is used throughout Linux, and \OSC\,\S1.9 covers the standard kernel data structures
in general.

\subsection{Testing a kernel automatically}

QEMU has a tiny virtual device made for testing, \code{isa-debug-exit}. When the guest writes a value $v$ to its I/O
port, QEMU exits immediately with status $2v + 1$. A kernel started with a \code{test} command line can run every
test, print the results on the serial port, and write 0 (exit status 1, meaning success) or 1 (status 3, failure) to
the port. A Makefile target turns that into an ordinary pass or fail, which you can run before every commit.

\section{Reading plan}

\begin{readingplan}
\must & \OSTEP & Ch.~17 \emph{Free-Space Management} (re-read with the implementation in mind: headers, free list, splitting, coalescing, strategies, segregated lists, buddy) & 161--178 & The design of Tasks 7.2--7.4. \\
\must & \OSTEP & Ch.~14 \emph{Interlude: Memory API} & 123--134 & Every bug in its list is something your heap must detect. \\
\must & \LOB & Ch.~10, section \emph{A Kernel Heap} & 63 & Short, and points to the same options. \\
\should & \OSC & \S10.8 \emph{Allocating Kernel Memory} (buddy system, slab allocation) & 426--430 & What production kernels do, and why. \\
\should & \OSC & \S1.9 \emph{Kernel Data Structures} & 36--40 & Lists, queues, trees, hash maps, bitmaps in kernels. \\
\could & \MOS & \S3.2.3 \emph{Managing Free Memory} & 190--194 & First, next, best, worst and quick fit compared. \\
\could & \MOS & \S1.8.3 \emph{Large Programming Projects} & 75--76 & Your source tree is growing; organise headers and objects now. \\
\end{readingplan}

\begin{minixbox}
The MINIX 3 kernel has no heap at all. Every table has a fixed size chosen at compile time: the process table is the
array \code{proc[NR_TASKS + NR_PROCS]} declared in \texttt{kernel/proc.h (05500)}, and \texttt{kernel/table.c
(06000)} holds other static tables. That makes memory use predictable and removes a whole class of allocator bugs, at
the price of hard limits on processes, open files and so on. Turk-OS takes the other road, but you could still use
fixed tables for some objects. Think about which one you would choose for the process table, and why.
\end{minixbox}

\section{Build it}

\task{Whole pages from the direct map}
Add \code{kalloc_page()} and \code{kfree_page()}. Because the direct map covers all usable RAM, a frame from the PMM
is already reachable at \code{P2V(frame)}, so these are two one-line functions. Zero pages on allocation while you
are debugging, so you never read stale data.

\begin{lst}{c}{kernel/mm/kalloc.c}
void *kalloc_page(void)
{
    uintptr_t pa = pmm_alloc_frame();
    if (!pa)
        return NULL;
    void *va = P2V(pa);
    memset(va, 0, PAGE_SIZE);
    return va;
}

void kfree_page(void *va)
{
    ASSERT(((uintptr_t)va & (PAGE_SIZE - 1)) == 0);
    pmm_free_frame(V2P(va));
}
\end{lst}

\task{A heap region that grows}
Reserve \code{KHEAP_START = 0xE0000000} and \code{KHEAP_MAX = 0xF0000000}. Keep a \code{heap_end} pointer, and write
\code{ksbrk(n)}, which maps enough fresh frames to extend the heap by at least $n$ bytes and returns the old end.
The page tables for this region already exist (Phase 6), so no kernel PDE ever changes.
\verify{Call \code{ksbrk(10000)} twice and check with \code{info mem} that three pages, then six, are mapped from
\code{0xE0000000}.}

\task{kmalloc and kfree}
Implement first fit with splitting and coalescing over a doubly linked list of blocks kept in address order. Make the
header a multiple of 16 bytes and round every request up to 16, so every pointer you return is 16-byte aligned. Put a
magic number in each header: it catches most buffer overflows and bad frees the moment they happen.

\begin{lst}{c}{kernel/mm/heap.c}
#define HEAP_MAGIC 0xC0FFEE42u
#define ALIGN16(x) (((x) + 15) & ~(size_t)15)

struct block {
    uint32_t      magic;
    uint32_t      size;        /* payload size in bytes, a multiple of 16  */
    uint32_t      free;
    uint32_t      pad;
    struct block *next;        /* neighbours in address order              */
    struct block *prev;
    uint32_t      pad2[2];
};                             /* 32 bytes: payload stays 16-byte aligned  */
_Static_assert(sizeof(struct block) % 16 == 0, "header must keep alignment");

void *kmalloc(size_t size)
{
    if (size == 0)
        return NULL;
    size = ALIGN16(size);
    for (;;) {
        for (struct block *b = heap_head; b; b = b->next) {
            if (b->free && b->size >= size) {
                split_block(b, size);          /* only if the rest can hold a header + 16 */
                b->free = 0;
                return b + 1;                  /* payload starts right after the header   */
            }
        }
        if (!grow_heap(size + sizeof(struct block)))
            return NULL;                       /* out of memory */
    }
}

void kfree(void *p)
{
    if (!p)
        return;
    struct block *b = (struct block *)p - 1;
    ASSERT(b->magic == HEAP_MAGIC);            /* bad pointer or corrupted header */
    ASSERT(!b->free);                          /* double free                     */
    b->free = 1;
    coalesce(b);                               /* merge with free neighbours      */
}
\end{lst}
\verify{Allocate three blocks, free the middle one, then the first: \code{heap_check()} (next task) finds one free
block covering both, not two.}

\task{The rest of the API, and a heap checker}
Add \code{kcalloc} (allocate and zero), \code{krealloc} (grow in place when the next block is free and large enough,
otherwise allocate, copy and free) and \code{kmalloc_stats()}. Then write \code{heap_check()}, which walks every block
and asserts the invariants: every magic is correct, \code{b->next->prev == b}, blocks are in address order and touch
each other exactly, and no two free blocks are adjacent. Add a monitor command \code{heap} that prints the block
count, bytes in use and free, and the largest free block.
\verify{\code{heap_check()} passes after each step of every test in Task~7.7.}

\task{An intrusive list}
Write \code{include/list.h} with a circular doubly linked list and a dummy head node, which removes every special
case for empty lists and list ends. Phase 8 builds the scheduler's queues on it, and many later structures follow.

\begin{lst}{c}{include/list.h}
#include <stddef.h>

struct list_node { struct list_node *next, *prev; };

#define container_of(ptr, type, member) \
    ((type *)((char *)(ptr) - offsetof(type, member)))

static inline void list_init(struct list_node *head) { head->next = head->prev = head; }
static inline int  list_empty(const struct list_node *head) { return head->next == head; }

static inline void list_add_tail(struct list_node *head, struct list_node *n)
{
    n->prev = head->prev;
    n->next = head;
    head->prev->next = n;
    head->prev = n;
}

static inline void list_del(struct list_node *n)
{
    n->prev->next = n->next;
    n->next->prev = n->prev;
    n->next = n->prev = NULL;      /* makes use-after-remove crash loudly */
}

#define list_for_each(pos, head) \
    for ((pos) = (head)->next; (pos) != (head); (pos) = (pos)->next)
\end{lst}
\verify{A test puts five structs on a list, removes the second and fourth, and walks it with \code{container_of}
to check the remaining values and their order.}

\task{make test}
Add \code{outl} to \code{io.h}. In \code{kmain}, look at the Multiboot command line: if it contains \code{test}, run
the whole suite, print a summary on serial, and exit QEMU through \code{isa-debug-exit}. Then add the Makefile target.
A timeout protects you from a test that hangs.

\begin{lst}{c}{kernel/kmain.c (test mode)}
static inline void qemu_exit(uint32_t v) { outl(0xF4, v); }   /* QEMU exits with 2v+1 */

if (cmdline_has(mbi, "test")) {
    int failed = ktest_run_all();          /* prints PASS/FAIL per test on serial */
    qemu_exit(failed ? 1 : 0);
}
\end{lst}

\begin{lst}{make}{Makefile (test target)}
QEMU_TEST := qemu-system-i386 -kernel $(BUILD)/kernel.elf -append "test" \
             -display none -serial stdio -no-reboot \
             -device isa-debug-exit,iobase=0xf4,iosize=0x04

test: $(BUILD)/kernel.elf
	@timeout 120 $(QEMU_TEST); status=$$?; \
	if [ $$status -eq 1 ]; then echo "ALL TESTS PASSED"; \
	else echo "TESTS FAILED (status $$status)"; exit 1; fi
\end{lst}
\verify{\code{make test} prints every test result and \code{ALL TESTS PASSED}. Break one test on purpose and
check that the target now fails with status 3.}

\task{Stress the heap}
Write a randomised test with a small pseudo-random generator (xorshift is enough). Keep 256 slots. At each of 10\,000
steps pick a slot: if it is empty, allocate a random size from 1 to 4096 bytes and fill it with the slot's number; if
it is full, check that the fill pattern is intact and free it. Run \code{heap_check()} every 100 steps. At the end,
free everything and check that bytes in use is zero and that the heap is one free block again.
\verify{The stress test passes, is deterministic for a given seed, and is part of \code{make test}.}

\task{A slab cache (optional)}
Add \code{kmem_cache_create(name, object_size)}, \code{kmem_cache_alloc} and \code{kmem_cache_free}: each cache takes
whole pages from \code{kalloc_page} and cuts them into equal objects linked on a free list. Use it for
\code{struct task} in Phase 8 and compare the heap statistics with and without it.

\section{Testing and debugging}

\begin{pitfalls}
\code{ASSERT(b->magic == HEAP_MAGIC)} fails in \code{kfree} & A buffer overflow in the previous block overwrote this header, or a pointer that didn't come from \code{kmalloc} was freed & Print the neighbouring headers; check the caller's sizes. Red zones (Going further) pinpoint the overflow. \\
Returned pointers are not 16-byte aligned & The header size, or a split size, is not a multiple of 16 & Keep the \code{_Static_assert}; round the split size up. \\
The heap keeps growing although usage is flat & Coalescing misses a case, so blocks fragment & \code{heap_check()} asserting ``no two adjacent free blocks'' finds it. \\
List corruption after removing a block & A \code{prev} pointer not updated in split or coalesce & Assert \code{b->next->prev == b} in \code{heap_check()}. \\
\code{make test} hangs forever & A test waits for a key, or the kernel crashed without reaching \code{qemu_exit} & The timeout catches it; make \code{panic} call \code{qemu_exit(1)} in test mode too. \\
\end{pitfalls}

\begin{warnbox}[A rule for later phases]
Don't call \code{kmalloc} or \code{kfree} from interrupt handlers. Once Phase 8 adds preemption and Phase 9 adds locks,
an interrupt that arrives while the heap lock is held would deadlock if its handler tried to take the same lock.
Interrupt handlers should use memory allocated in advance.
\end{warnbox}

\section{Milestone}

\begin{milestonebox}[The kernel can allocate]
\begin{checklist}
  \item \code{kalloc_page}/\code{kfree_page} and \code{kmalloc}/\code{kfree}/\code{kcalloc}/\code{krealloc} work, with 16-byte alignment.
  \item Magic numbers catch double frees and bad pointers; \code{heap_check()} verifies every invariant.
  \item The 10\,000-step stress test passes, and all memory returns to the heap at the end.
  \item \code{list.h} is tested and ready for the scheduler.
  \item \code{make test} runs headless and fails loudly when any test fails.
\end{checklist}
\end{milestonebox}

\begin{questionbox}
\begin{enumerate}
  \item How does \code{kfree(p)} know how big the block is?
  \item What is coalescing, when does it happen, and what goes wrong without it?
  \item First fit or best fit: which fragments memory faster, and which searches faster?
  \item Why can't \code{kmalloc} hand out whole pages for page tables? What does Turk-OS use instead?
  \item How does a slab allocator reduce both fragmentation and initialisation cost?
  \item What exactly does \code{container_of} compute? Write it out for a \code{struct task} whose list node is at offset 24.
\end{enumerate}
\end{questionbox}

\section{Going further}

Add a debug mode with \textbf{red zones}: a few guard bytes with a known value after each allocation, checked on
\code{kfree}, which pinpoints overflows. Fill freed memory with \code{0xDEADBEEF} so use-after-free shows up as
obviously wrong values. Record the caller of every allocation with \code{__builtin_return_address(0)} and add a
\code{leaks} command that lists live allocations by caller. Or replace first fit with segregated free lists (one list
per size class), and measure the difference with the stress test.

\nextphase{8}{Processes, Context Switching and Scheduling}
\booklegend
\end{document}
